\section{Parameter Choices and Deviations}
\label{sec:deviations}

For SFT and both DPO runs, I kept the guide recipe wherever nothing forced a change. Every headline hyperparameter in Table~\ref{tab:config} matches the guide tables, and so does the 256-token DPO prompt limit in Table~\ref{tab:data}. Table~\ref{tab:deviations} lists every other change, grouped by whether resource limits or performance evidence drove it.

\begin{table}[!t]
\centering
\caption{Changes from the guide, grouped by what drove them. ``My choice'' means the guide leaves the setting to me.}
\label{tab:deviations}
\begingroup\tablefont\setlength{\tabcolsep}{3pt}\renewcommand{\arraystretch}{1.0}
\begin{tabularx}{\textwidth}{>{\raggedright\arraybackslash}p{2.3cm}>{\raggedright\arraybackslash}p{2.5cm}>{\raggedright\arraybackslash}p{3.1cm}>{\raggedright\arraybackslash}X}
\toprule
Setting & Guide & Mine & Reason and effect \\
\midrule
\rowcolor{accent!8}\multicolumn{4}{@{}l}{\textbf{Resource-driven}} \\
SFT validation batch & not specified & 4, down from 8 & Validation ran out of memory on the T4. The loss is an exact token mean, so no value changes. \\
\rowcolor{accent!8}\multicolumn{4}{@{}l}{\textbf{Data and implementation choices}} \\
Long SFT examples & truncate to the maximum length & skip rows over 512 tokens until 30,000 fit & Truncation would cut the EOS label, so the model would never learn to stop. I did not test it. \\
DPO reference precision & fp16 allowed & fp32 weights under fp16 autocast & Memory allowed fp32, so policy and reference start identical. \\
\rowcolor{accent!8}\multicolumn{4}{@{}l}{\textbf{Performance-driven (LoRA persona task, where the guide leaves these settings to me)}} \\
Loss & plain DPO & DPO plus NLL on the chosen answer, weight 1 & Plain DPO left no persona (style 0/40/0). NLL gave small, mixed changes. \\
Learning rate & my choice & $5\times10^{-5}$ & The margin barely moved at $1\times10^{-5}$, and EOS collapsed from $7\times10^{-5}$ up (Table~\ref{tab:lorasweep}). \\
Epochs & my choice & 3 & Ten epochs raised style but cut EOS to 11/40, and the filter rejected it. \\
Rank & my choice & 16 (alpha 32) & Rank 8 reached half the margin at $5\times10^{-5}$ (0.52 against 1.02). I did not try it at a higher rate. \\
Targets and batch & my choice & all seven linear projections; 4 $\times$ 4 = 16 & Not ablated. A batch of 16 gives 13 updates per epoch, and 32 would give 7. \\
\bottomrule
\end{tabularx}
\endgroup

\end{table}

\paragraph{Resource-driven changes.} The memory problem I hit was in SFT validation, not in training. Validation ran out of memory on the T4, where training itself peaked at 12.8 GB, so I halved the validation batch from 8 to 4 and set the expandable-segments allocator. Validation loss is an exact token mean, so the batch size changes no value. Memory never forced a training change: SFT and DPO ran at the guide batch sizes without gradient checkpointing.

\paragraph{Data and implementation choices.} Three choices differ from the guide without anything forcing them. The guide truncates long SFT examples to 512 tokens, and I skipped them instead (21,560 of the 51,565 rows I read) until 30,000 fit, because truncation would cut the EOS label and the model would never learn to stop. The guide lets the DPO reference load in fp16, and I used fp32: it needed only 3.4 to 4.2 GB on its own T4, and policy and reference then start identical, with the largest absolute implicit reward at step 0 at $2.4\times10^{-5}$ for HH-RLHF and $4.3\times10^{-5}$ for UltraFeedback.

\paragraph{LoRA settings.} The guide leaves the LoRA settings to me, so these are the performance-driven choices, and I made them by training 12 adapters in two sweeps (Table~\ref{tab:lorasweep}). Plain DPO at $5\times10^{-5}$ gave no persona on the 40 required prompts (0 wins, 40 ties, 0 losses), so I added the NLL term on the chosen answer. It changed little: held-out EOS went from 19 to 20 of 22, and required-prompt EOS from 40 to 35 of 40. At $1\times10^{-5}$ the validation margin barely moved (0.12), and from $7\times10^{-5}$ up the answers degraded: required-prompt EOS fell from 40 of 40 at $5\times10^{-5}$ to 31, 30, and 2 of 40 at $7\times10^{-5}$, $1\times10^{-4}$, and $2\times10^{-4}$. More epochs did the same. Ten epochs at $5\times10^{-5}$ gave the strongest plain-DPO style (9 wins, 31 ties, 0 losses), but EOS fell to 11 of 40, and my filter rejected it. I chose rank 16 over rank 8 because rank 8 reached only half the margin at $5\times10^{-5}$ (0.52 against 1.02), though I did not try it at a higher rate. I did not ablate the target modules or the batch size (Table~\ref{tab:deviations}).

\paragraph{Which adapters I scored.} The dashes in Table~\ref{tab:lorasweep} show which checks each adapter went through. In the first sweep, I trained seven plain-DPO adapters for three epochs and generated answers only for the 40 required prompts, which gives their EOS and loop counts. I judged style for two of them (the $5\times10^{-5}$ and $1\times10^{-4}$ runs) and generated held-out answers only for the $5\times10^{-5}$ run. The second sweep added five adapters with more epochs or the NLL term, and I scored each on both prompt sets and judged its style. I did not take the other six first-sweep adapters further. Four of them, at learning rates of $7\times10^{-5}$ and above, already break the filter's thresholds on the 40 required prompts alone, which the table marks as ``fail on the 40 alone''. The other two, at $1\times10^{-5}$ and with rank 8, trained the smallest margins (0.12 and 0.52).

\paragraph{Adapter selection.} Before I scored any adapter on the held-out prompts, I fixed the selection rule. An adapter fails if, on either prompt set, its EOS rate falls by more than 10 points or its repeated-phrase rate rises by more than 10 points relative to the base model. Among the adapters that pass, I take the highest held-out style score. Only one of the six adapters I scored on both prompt sets passed, so I selected rank 16, a learning rate of $5\times10^{-5}$, three epochs, and the NLL term. That choice has a cost: the adapters that move the persona style most all fail the filter, and the one that passes barely moves it. Section~\ref{sec:results} reports both.

\begin{table}[!t]
\centering
\caption{All 12 adapters I trained for the persona task, in two sweeps. The columns marked 40 and the style column use the 40 required prompts, and the columns marked 22 use the 22 held-out prompts. EOS counts answers that end on their own, loops counts answers with a repeated phrase, and Style W/T/L is the judge's persona-style win, tie, and loss count against the base model. A dash means I did not run that check. The filter needs both prompt sets, so ``fail on the 40 alone'' marks adapters that already break its 10-point thresholds on the required prompts. Shading marks the selected adapter.}
\label{tab:lorasweep}
\begingroup\tablefont\setlength{\tabcolsep}{3pt}\renewcommand{\arraystretch}{1.0}
\begin{tabularx}{\textwidth}{>{\centering\arraybackslash}p{0.8cm}>{\centering\arraybackslash}p{1.4cm}>{\centering\arraybackslash}p{1.0cm}>{\centering\arraybackslash}p{1.9cm}>{\centering\arraybackslash}p{1.2cm}>{\centering\arraybackslash}p{1.2cm}>{\centering\arraybackslash}p{1.2cm}>{\centering\arraybackslash}p{1.2cm}>{\centering\arraybackslash}p{1.2cm}>{\centering\arraybackslash}p{1.7cm}>{\raggedright\arraybackslash}X}
\toprule
Rank & Learning rate & Epochs & Loss & Val. margin & EOS of 40 & Loops of 40 & EOS of 22 & Loops of 22 & Style W/T/L & Filter \\
\midrule
\multicolumn{5}{@{}l}{Base model, no adapter} & 37/40 & 11 & 22/22 & 2 & - & - \\
16 & $1\times10^{-5}$ & 3 & DPO & 0.12 & 38/40 & 12 & - & - & - & - \\
8 & $5\times10^{-5}$ & 3 & DPO & 0.52 & 39/40 & 11 & - & - & - & - \\
16 & $5\times10^{-5}$ & 3 & DPO & 1.02 & 40/40 & 11 & 19/22 & 5 & 0/40/0 & fail \\
16 & $7\times10^{-5}$ & 3 & DPO & 1.55 & 31/40 & 15 & - & - & - & fail on the 40 alone \\
16 & $1\times10^{-4}$ & 3 & DPO, $\beta=0.3$ & 5.17 & 33/40 & 20 & - & - & - & fail on the 40 alone \\
16 & $1\times10^{-4}$ & 3 & DPO & 2.30 & 30/40 & 24 & - & - & 2/38/0 & fail on the 40 alone \\
16 & $2\times10^{-4}$ & 3 & DPO & 5.47 & 2/40 & 38 & - & - & - & fail on the 40 alone \\
16 & $5\times10^{-5}$ & 10 & DPO & 4.79 & 11/40 & 33 & 3/22 & 20 & 9/31/0 & fail \\
\rowcolor{accent!10}16 & $5\times10^{-5}$ & 3 & DPO + NLL & 0.96 & 35/40 & 13 & 20/22 & 3 & 0/40/0 & pass (selected) \\
16 & $1\times10^{-4}$ & 3 & DPO + NLL & 2.07 & 28/40 & 18 & 10/22 & 12 & 1/39/0 & fail \\
16 & $5\times10^{-5}$ & 6 & DPO + NLL & 2.21 & 32/40 & 17 & 10/22 & 12 & 1/39/0 & fail \\
16 & $1\times10^{-4}$ & 6 & DPO + NLL & 3.59 & 33/40 & 12 & 15/22 & 9 & 24/16/0 & fail \\
\bottomrule
\end{tabularx}
\endgroup

\end{table}
